\setcounter{chapter}{9}
\renewcommand{\thetheorem}{10.a.\arabic{theorem}}
\setcounter{theorem}{0}

\chapter{Bases of Vector Spaces (Part A)}

\section{Foundations of Bases and the Exchange Theorem}

\subsection{Definition and Unique Representation}

We have previously explored two fundamental properties a set of vectors can possess: the ability to \textit{span} the entire space, and the property of being \textit{linearly independent} (having no redundancies). A basis is simply a set that achieves both simultaneously.

\begin{definition}[Basis]
Let $V$ be a vector space over a field $K$. A subset $S \subseteq V$ is called a \textbf{basis} of $V$ (plural: \qt{bases}) if the following two conditions hold:
\begin{enumerate}
    \item $S$ is linearly independent.
    \item $S$ generates (or spans) $V$, meaning $\Sp(S) = V$.
\end{enumerate}
\end{definition}

Why is a basis so mathematically desirable? Because it provides a perfect "coordinate system" for our abstract vector space.

\begin{proposition}[Unique Representation]
A subset $S \subseteq V$ is a basis of $V$ if and only if every vector $v \in V$ can be written in a \textbf{unique way} as a linear combination of vectors in $S$.
\end{proposition}

To prove this rigorously—especially since $S$ might be an infinite set, like the set of all polynomials—we must first formalize what we mean by "a unique way as a linear combination." In a vector space, we can only ever add a \textit{finite} number of vectors together. 

We formalize this using \textbf{finitely supported functions}. Let $c: S \to K$ be a function that assigns a scalar coefficient to every vector in $S$. We require that $c$ is non-zero on only finitely many elements. In other words, the "support" of the function:
\[
\text{supp}(c) := \{u \in S \mid c(u) \ne 0\}
\]
must be a finite set. We can then define our linear combination as the sum:
\[
\sum_{u \in S} c(u) \cdot u
\]
Because $\text{supp}(c)$ is finite, the number of non-zero terms in this sum is strictly finite, making it a perfectly valid vector operation. (Furthermore, because vector addition is commutative, the order in which we sum these elements does not matter).

\begin{center}
\begin{tikzpicture}[
    node distance=1.5cm,
    box/.style={draw, thick, rounded corners, fill=blue!5, align=center, inner sep=2ex},
    arrow/.style={->, thick, >=stealth, draw=blue!60!black}
]
    \node[box] (funcs) {\textbf{Finitely Supported Functions} \\ $c: S \to K$ \\ (Only finitely many $c(u) \ne 0$)};
    \node[box, right=3cm of funcs] (vectors) {\textbf{Vectors in $V$} \\ $v = \sum_{u \in S} c(u) u$};
    
    \draw[arrow] (funcs.east) to[bend left=15] node[midway, above] {Linear Combination} (vectors.west);
    \draw[arrow, dashed] (vectors.west) to[bend left=15] node[midway, below] {Unique Recipe (if Basis)} (funcs.east);
\end{tikzpicture}
\end{center}

To say that every $v \in V$ can be written in a unique way means: if we have two finitely supported functions $c, c': S \to K$ such that:
\[
\sum_{u \in S} c(u) \cdot u = \sum_{u \in S} c'(u) \cdot u
\]
then it must be true that $c(u) = c'(u)$ for all $u \in S$. 

\begin{proof}[Proof of Proposition 10.1.1]
We must prove both directions.

\textbf{Direction 1 ($\implies$):} Assume $S$ is a basis. Because $S$ spans $V$, every $v \in V$ can be written as a linear combination of elements in $S$. We must show this combination is unique. 
Suppose there are two combinations yielding the same vector: $\sum_{u \in S} c(u) u = \sum_{u \in S} c'(u) u$. 
Subtracting one from the other yields:
\[
\sum_{u \in S} \big(c(u) - c'(u)\big) u = 0_V
\]
Because $S$ is a basis, it is linearly independent. The only way a linear combination of independent vectors can equal $0_V$ is if all coefficients are zero. Thus, $c(u) - c'(u) = 0 \implies c(u) = c'(u)$ for all $u \in S$. The representation is unique.

\textbf{Direction 2 ($\Longleftarrow$):} Assume every $v \in V$ has a unique representation. Since every vector can be represented, $S$ clearly spans $V$. We must show $S$ is independent. 
Consider the linear combination $\sum_{u \in S} c(u) u = 0_V$. We already know one valid representation for the zero vector: assigning $0$ to every coefficient. Since the hypothesis guarantees representations are \textit{unique}, the all-zero assignment must be the \textit{only} representation. Thus, $S$ is linearly independent. 
\end{proof}

\subsection{Standard Examples of Bases}

Let us identify the standard bases for the families of vector spaces we encounter most frequently.

\begin{example}[The Coordinate Space]
In $V = K^n$, the standard basis is $S = \{e_1, e_2, \dots, e_n\}$, where $e_i$ is the column vector with a $1$ in the $i$-th position and $0$ elsewhere.
\end{example}

\begin{example}[Bounded Polynomials]
In $V = K[x]_d$, the space of polynomials of degree $\le d$, the standard basis is $S = \{1, x, x^2, \dots, x^d\}$. 
\end{example}

\begin{example}[All Polynomials]
In $V = K[x]$, the space of all polynomials, the standard basis is the infinite set $S = \{1, x, x^2, \dots, x^n, \dots\}$. (Notice how beautifully our "finitely supported function" handles this: any specific polynomial has a finite degree, so it requires only a finite number of non-zero coefficients from this infinite basis).
\end{example}

\begin{example}[Matrices]
In $V = \M_{m \times n}(K)$, the standard basis is $S = \{E_{i,j} \mid 1 \le i \le m, 1 \le j \le n\}$, where $E_{i,j}$ is the matrix with a $1$ in the $(i,j)$ entry and $0$ everywhere else.
\end{example}

\subsection{The Size of a Basis and Exchange Lemmas}

Looking at $K^n$, the standard basis has exactly $n$ elements. Looking at $K[x]_d$, the standard basis has $d+1$ elements. This raises a monumental question:
\textbf{Are all bases of a given vector space $V$ exactly the same size?}

Yes, they are. Our next major goal is to prove this fact. To do so, we must understand how vectors can be swapped in and out of spanning sets without breaking the span. We develop this through a sequence of four powerful lemmas.

\begin{lemma}[Lemma A: Simple Exchange]
Let $v_1, \dots, v_n$ be a basis of $V$. Let $u = a_1 v_1 + \dots + a_n v_n \in V$ be a vector such that $a_1 \ne 0$. Then the modified list $u, v_2, \dots, v_n$ is also a basis of $V$.
\end{lemma}

\begin{proof}
We must show the new list is independent and spans $V$.
\textbf{Independence:} Assume $b_1 u + b_2 v_2 + \dots + b_n v_n = 0_V$. Substitute the definition of $u$:
\[
b_1 (a_1 v_1 + a_2 v_2 + \dots + a_n v_n) + b_2 v_2 + \dots + b_n v_n = 0_V
\]
Regrouping the terms by the basis vectors $v_i$:
\[
(b_1 a_1) v_1 + (b_1 a_2 + b_2) v_2 + \dots + (b_1 a_n + b_n) v_n = 0_V
\]
Since $v_1, \dots, v_n$ are linearly independent, all coefficients must be zero. In particular, $b_1 a_1 = 0$. Because we assumed $a_1 \ne 0$, we are forced to conclude $b_1 = 0$. Once $b_1 = 0$, the equation reduces to $b_2 v_2 + \dots + b_n v_n = 0_V$, which forces $b_2 = \dots = b_n = 0$. Thus, the new list is independent.

\textbf{Spanning:} We can rearrange $u = a_1 v_1 + a_2 v_2 + \dots + a_n v_n$ to solve for $v_1$:
\[
v_1 = a_1^{-1} u - a_1^{-1} a_2 v_2 - \dots - a_1^{-1} a_n v_n
\]
This shows $v_1 \in \Sp(u, v_2, \dots, v_n)$. Since all the other basis vectors $v_2, \dots, v_n$ are trivially in this span, all vectors of the original basis are contained within it. Thus $V = \Sp(v_1, \dots, v_n) \subseteq \Sp(u, v_2, \dots, v_n)$.
\end{proof}

\begin{lemma}[Lemma B: Targeted Redundancy]
Let $v_1, \dots, v_n$ be a list of vectors. Suppose it is linearly dependent. Then there exists an index $1 \le j \le n$ such that $v_j \in \Sp(v_1, \dots, v_{j-1})$. Moreover, if we drop $v_j$ from the list, the span of the remaining vectors is exactly the same as the original span.
\end{lemma}

\begin{proof}
Because the list is dependent, there exist coefficients not all zero such that $\sum_{i=1}^n a_i v_i = 0_V$. Let $j$ be the \textit{largest} index such that $a_j \ne 0$. Then our equation is actually:
\[
a_1 v_1 + a_2 v_2 + \dots + a_j v_j = 0_V
\]
Because $a_j \ne 0$, we can solve for $v_j$:
\[
v_j = -\frac{a_1}{a_j} v_1 - \dots - \frac{a_{j-1}}{a_j} v_{j-1}
\]
This proves $v_j \in \Sp(v_1, \dots, v_{j-1})$. Because $v_j$ can be synthesized entirely from the vectors strictly preceding it, dropping $v_j$ does not shrink the span.
\end{proof}

\begin{lemma}[Lemma C: Span Overcrowding]
If $\Sp(v_1, \dots, v_n) = V$ and $u \in V$, then the list $u, v_1, \dots, v_n$ is linearly dependent.
\end{lemma}

\begin{proof}
Because the $v_i$'s span $V$, we can write $u = a_1 v_1 + \dots + a_n v_n$. Rearranging this gives:
\[
1 \cdot u - a_1 v_1 - \dots - a_n v_n = 0_V
\]
Since the coefficient of $u$ is $1$ (which is non-zero), this is a non-trivial linear combination that equals zero. The list is dependent.
\end{proof}

We now arrive at the crescendo of this chapter: The Steinitz Exchange Lemma. This theorem rigorously proves that an independent set can never outnumber a spanning set.

\begin{theorem}[Lemma D: Steinitz Exchange Theorem]
Let $u_1, \dots, u_m$ be a linearly independent list in $V$. Let $v_1, \dots, v_n$ be a list that spans $V$ (i.e., $\Sp(v_1, \dots, v_n) = V$). Then $m \le n$.
\end{theorem}

\begin{center}
\begin{tikzpicture}[
    box/.style={draw, thick, rounded corners, align=center, inner sep=2ex},
    arrow/.style={->, thick, >=stealth, draw=red!80!black}
]
    \node[box, fill=green!10] (indep) {\textbf{Independent List} ($m$ elements) \\ $u_1, u_2, \dots, u_m$};
    \node[box, fill=blue!10, right=3cm of indep] (span) {\textbf{Spanning List} ($n$ elements) \\ $v_1, v_2, \dots, v_n$};
    
    \draw[arrow] (indep.north east) to[bend left=30] node[midway, above] {\textbf{Step 1:} Insert $u_1$} (span.north west);
    
    \node[box, fill=white, dashed, below=1cm of span] (overflow) {List becomes dependent (Lemma C). \\ A redundant $v_k$ is ejected (Lemma B). \\ \textbf{Length remains exactly $n$.}};
    
    \draw[->, thick, dashed] (span.south) -- (overflow.north);
\end{tikzpicture}
\end{center}

\begin{proof}
We proceed by injecting the independent $u$ vectors into the spanning list one by one, forcefully replacing $v$ vectors to maintain a constant list length of $n$.

\textbf{Step 1:} Add $u_1$ to the front of the spanning list: $u_1, v_1, \dots, v_n$. Because the $v_i$'s span $V$, Lemma C dictates this new list is dependent. By Lemma B, there is a redundant vector in this list that can be dropped without altering the span. 
Could the redundant vector be $u_1$? No. Lemma B specifies the redundant vector is in the span of the vectors \textit{preceding} it. $u_1$ is at the very front (its preceding span is $\emptyset$), and $u_1 \ne 0_V$ because it is part of an independent set. Thus, the redundant vector must be one of the $v_i$'s, say $v_k$. We drop $v_k$. 
Our list is now $u_1, v_1, \dots, \widehat{v_k}, \dots, v_n$. It still spans $V$, and it has exactly $n$ elements.

\textbf{Step 2:} Add $u_2$ to the front: $u_2, u_1, \dots, \widehat{v_k}, \dots, v_n$. By Lemma C, this is dependent. By Lemma B, a redundant vector exists. It cannot be $u_2$ (front of the line). It cannot be $u_1$, because $u_1$ cannot be in the span of $u_2$ (since the $u$'s are independent). Therefore, the redundant vector \textit{must} be another $v$. We drop it. Length remains $n$.

\textbf{The Contradiction:} Assume, for the sake of contradiction, that $m > n$. 
We can continue the above process exactly $n$ times. After step $n$, we will have systematically ejected every single $v_i$ from the list, leaving us with a list consisting purely of $u$-vectors: $u_n, u_{n-1}, \dots, u_1$. 
Crucially, because we only ever used Lemma B to drop vectors, this list of $u$-vectors \textbf{still spans $V$}. 

Because $m > n$, we still have at least one vector left in our independent set: $u_{n+1}$. Let us add it to the front of our current spanning list:
\[
u_{n+1}, u_n, u_{n-1}, \dots, u_1
\]
Because the set $u_n, \dots, u_1$ spans $V$, Lemma C states that this new combined list must be linearly dependent. However, this list is just a subset of our original $u_1, \dots, u_m$, which was explicitly assumed to be linearly independent! This is a fatal mathematical contradiction.

Therefore, our assumption that $m > n$ must be false. We conclude that $m \le n$.
\end{proof}
